% Blender primarily uses the Mantaflow fluid simulation framework, which replaced older methods and supports both liquid and gas simulations. The specific algorithm Mantaflow uses is a hybrid approach, but older versions of Blender relied on the Lattice Boltzmann Method (LBM) for liquids, and some effects can be created using Smoothed Particle Hydrodynamics (SPH). 


% Численное решение уравнений движения (здесь и дальше \textemdash{} уравнений Навье-Стокса)
% сводится к последовательному обновлению значений этих полей по шагам времени.
% Такой подход хорошо согласуется с архитектурой графических процессоров:
% каждая ячейка сетки может обрабатываться независимым потоком, что позволяет
% эффективно реализовывать эйлеровские методы с использованием WebGPU.


В компьютерной графике для моделирования течения жидкости
широко применяются различные численные методы. Выбор конкретного алгоритма
зависит от требований к качеству визуализации, производительности и
сложности сцены.

Как правило, от выбранного алгоритма ожидаются:

\begin{itemize}
      \item поддержка реалистичных, но в то же время
            визуально выразительных и скорее красивых эффектов;
      \item возможность эффективной реализации на GPU
            в виде набора вычислительных шейдеров;
      \item устойчивость при достаточно крупных шагах по времени;
      \item разумная сложность реализации;
\end{itemize}

Рассмотрим несколько современных подходов,
широко применяемых для моделирования и визуализации
жидкостей.

\subsection{Метод Stable Fluids}

Метод Stable Fluids \cite{stamstablefluids} относится к эйлеровым,
сеточным методам решения уравнений Навье-Стокса для
несжимаемой жидкости. Пространство дискретизуется равномерной
прямоугольной сеткой; в ячейках (или на их границах) хранятся
поле скорости, давление и дополнительные скалярные поля
(плотность дыма, температура, цвет и т.п.). Подход был
предложен Й.~Стамом специально для задач интерактивной
компьютерной графики и основан на неявной, безусловно устойчивой
схеме интегрирования по времени.

Исходной точкой служат уравнения Навье-Стокса для
несжимаемой жидкости \eqref{eq:NS}.~ Используя теорему о разложении Гельмгольца, векторное поле скорости $\mathbf{w}$ можно единственным образом разложить в виде:
\begin{equation}\label{eq:helmholtz}
      \mathbf{w} = \mathbf{u} + \nabla q,
\end{equation}
где $\nabla\cdot\mathbf{u} = 0$, а $q$~ \textemdash{} скалярное поле \cite{stamstablefluids}. Такое разложение позволяет выделить бездивергентную (соленоидальную) часть поля и градиентную составляющую. Это приводит к возможности введения оператора $P$, который вычитает градиент давления
и делает поле скорости несжимаемым. В компактном виде
уравнение скорости записывается как
\begin{equation}\label{eq:NS_projection}
      \frac{\partial \mathbf{u}}{\partial t}
      =
      P\!\left(
      -(\mathbf{u}\cdot\nabla)\mathbf{u}
      + \nu \nabla^2 \mathbf{u}
      + \mathbf{f}
      \right).
\end{equation}

В численном алгоритме Stable Fluids один шаг по времени
разбивается на несколько подшагов (operator splitting). Пусть
$\mathbf{u}^n$ \textemdash{} поле скорости в момент времени $t^n$. Тогда
поле $\mathbf{u}^{n+1}$ для времени $t^{n+1}=t^n+\Delta t$
получается последовательным применением следующих стадий:

\begin{enumerate}
      \item \textbf{Добавление внешних сил.}
            На этом шаге учитываются внешние воздействия такие как гравитация и прочее.
            Предполагая, что $\mathbf{f}$ мало меняется за интервал $\Delta t$, получаем
            промежуточное поле
            \begin{equation}
                  \mathbf{w}_1
                  =
                  \mathbf{u}^n + \Delta t\,\mathbf{f}.
            \end{equation}

      \item \textbf{Адвекция (полулагранжева схема).}
            Ключевая идея Stable Fluids \textemdash{} использовать полулагранжевую
            адвекцию, основанную на методе характеристик. Для каждого
            центра ячейки в точке $\mathbf{x}$ трассируется обратная
            траектория на время $\Delta t$ назад по полю скорости
            $\mathbf{w}_1$:
            \begin{equation}
                  \mathbf{x}_0
                  =
                  \mathbf{x}
                  - \Delta t\,\mathbf{w}_1(\mathbf{x}).
            \end{equation}
            Новое значение скорости в точке $\mathbf{x}$ берётся как
            интерполированное значение старого поля в точке $\mathbf{x}_0$:
            \begin{equation}
                  \mathbf{w}_2(\mathbf{x})
                  =
                  \mathbf{w}_1(\mathbf{x}_0).
            \end{equation}
            В дискретном виде это означает: для каждой ячейки выполняется
            обратный шаг частицы и билинейная (в 2D) или трилинейная (в
            3D) интерполяция по соседним ячейкам. Такая схема
            \emph{безусловно устойчива} и допускает довольно крупные шаги
            по времени, хотя и вводит численную диффузию.

      \item \textbf{Диффузия (вязкость).}
            Вязкость описывается диффузионным слагаемым
            $\nu\nabla^2 \mathbf{u}$. В Stable Fluids он интегрируется
            неявно, что также даёт безусловную устойчивость по
            отношению к параметру $\nu$. Дискретизация по времени приводит
            к уравнению вида
            \[
                  \bigl(I - \nu \Delta t \nabla^2\bigr)\,\mathbf{w}_3
                  =
                  \mathbf{w}_2,
            \]
            где $I$ \textemdash{} тождественный оператор. После дискретизации
            лапласиана по сетке это уравнение превращается в разреженную
            систему линейных алгебраических уравнений, которая решается
            итерационными методами (Якоби, Гаусса\textemdash{}Зейделя, многосеточные
            методы и т.д.).

      \item \textbf{Проекция на несжимаемое поле (решение для давления).}
            После стадий адвекции и диффузии поле $\mathbf{w}_3$ вообще
            говоря не удовлетворяет условию $\nabla\cdot\mathbf{u}=0$.
            Чтобы восстановить несжимаемость, решается уравнение Пуассона
            для скалярного поля давления:
            \[
                  \nabla^2 p
                  = \rho\,\nabla\cdot\mathbf{w}_3,
            \]
            после чего из скорости вычитается градиент давления:
            \[
                  \mathbf{u}^{n+1}
                  =
                  \mathbf{w}_3
                  - \frac{\Delta t}{\rho}\,\nabla p.
            \]
            В дискретной форме это снова задача решения разреженной
            линейной системы. В случае периодических граничных условий
            возможно особенно эффективное решение через БПФ (FFT), так
            как лапласиан диагонализуется в фурье\textemdash{}базисе.
\end{enumerate}

Для дополнительных скалярных полей (плотность дыма, температура,
текстурные координаты) используется аналогичный шаблон:
источник $\to$ адвекция по уже вычисленному полю скорости
$\mathbf{u}^{n+1}$ $\to$ диффузия и (при необходимости)
диссипация. С точки зрения реализации на GPU каждая стадия
естественно оформляется отдельным проходом или compute\textemdash{}шейдером
над регулярной текстурой, что делает метод удобным для WebGPU
и интерактивных приложений в браузере.

Для метода Stable Fluids можно выделить следующие преимущества:

\begin{itemize}
      \item безусловная численная устойчивость, позволяющая использовать
            относительно крупные шаги по времени без дестабилизации решения;
      \item простая регулярная структура данных, хорошо подходящая
            для реализации на GPU (двумерные и трёхмерные текстуры,
            вычислительные шейдеры);
      \item естественная поддержка взаимодействия с внешними силами
            и дополнительными скалярными полями (дым, температура,
            текстурные координаты);
      \item хорошая визуальная правдоподобность для дыма, газа и
            двумерных симуляций воды при сравнительно низком разрешении
            сетки.
\end{itemize}

Однако данный подход имеет и свои ограничения. Полулагранжева
адвекция и неявная диффузия вносят заметную численную диссипацию (рассеивание / dissipation),
которая сглаживает мелкие вихри и детали и приводит к потере
вихревости (vorticity). Для компенсации этого эффекта часто
используют дополнительные приёмы вроде vorticity confinement.
Строгая привязка к равномерной сетке затрудняет адаптивное
увеличение разрешения в некоторых локальных областях. В
трёхмерном случае же вычислительная сложность возрастает
квадратично по числу ячеек в одном измерении, что делает реализацию
полноценного 3D Stable Fluids слишком ресурсоёмкой.

\subsection{SPH (Smoothed Particle Hydrodynamics)}

Метод SPH (Smoothed Particle Hydrodynamics) \cite{chinapaper}
относится к лагранжевым частичным методам. Жидкость представляется
набором частиц с массой, положением и скоростью, которые
перемещаются вместе с потоком. Каждая частица интерпретируется
как маленький объём жидкости, в котором усреднены плотность,
давление и другие величины.

Ключевая идея SPH \textemdash{} заменить непрерывные поля (плотность,
давление, скорость) на их приближенные (сглаженные) оценочные значения, вычисляемыми по
соседним частицам. Для этого вводится сглаживающее ядро
$W(\mathbf{r}, h)$ с радиусом влияния $h$. Плотность в точке,
соответствующей частице $i$, восстанавливается суммированием
влияния соседей:
\[
      \rho_i \approx \sum_{j} m_j \,
      W(\mathbf{x}_i - \mathbf{x}_j, h),
\]
где $m_j$ \textemdash{} масса частицы $j$, $\mathbf{x}_j$ \textemdash{} её положение.
Давление обычно задаётся через уравнение состояния
$p = p(\rho)$, а силы давления, вязкости и внешние силы
получаются из дискретизации уравнений Навье\textemdash{}Стокса
в частично\textemdash{}ориентированной форме.

Численный шаг SPH обычно включает следующие этапы:
обновление плотности по соседям, вычисление давления,
подсчёт сил (давление, вязкость, внешние поля, взаимодействие
с границами) и интегрирование уравнений движения для
обновления скоростей и положений частиц. Для ускорения
поиска соседей вводятся пространственные структуры данных
(регулярная решетка, хеш\textemdash{}таблица, деревья и т.п.), так как
каждая частица взаимодействует только с соседями в радиусе $h$.

Для метода SPH можно выделить следующие преимущества:

\begin{itemize}
      \item естественная поддержка свободной поверхности,
            брызг и разрывов струи, так как частицы просто
            разлетаются в пространстве;
      \item удобная работа со сложной геометрией и подвижными
            твёрдыми телами: границы задаются через специальные
            граничные частицы или коллизии;
      \item хорошо расширяется на 3D, так как алгоритм
            естественно трёхмерен и не привязан к структурированной
            сетке;
      \item лагранжева природа алгоритма упрощает адвекцию: не нужно
            явно решать задачу переноса на эйлеровой сетке;
\end{itemize}

А также следующие недостатки:

\begin{itemize}
      \item необходимость эффективного поиска соседей (дополнительные
            структуры данных и накладные расходы по памяти);
      \item повышенная вычислительная стоимость для плотных сцен
            и больших 2D/3D облаков частиц;
      \item настройка параметров сглаживания (радиус $h$, форма
            ядра, коэффициенты в уравнении состояния и вязкости)
            сильно влияет на стабильность и артефакты;
      \item сложность обеспечения строгой несжимаемости,
            требующая специальных вариаций SPH (WCSPH, PCISPH,
            DFSPH и др.);
\end{itemize}

В контексте визуализации течения жидкости для браузера SPH
представляет собой привлекательный вариант благодаря простой природе частиц
и наглядности, однако требует непростой реализации поиска соседей и выбора параметров, чтобы
добиться устойчивой работы в реальном времени.

\subsection{Модель мелкой воды (Shallow Water)}

Модель мелкой воды описывает движение тонкого слоя
жидкости через уравнения для высоты свободной
поверхности и средних по глубине скоростей. В простейшем
случае используется сетка высот, на которой решаются
гиперболические уравнения баланса.

Для этой модели можно выделить следующие преимущества:

\begin{itemize}
      \item высокая вычислительная эффективность;
      \item естественная визуализация волн на поверхности
            (например, вода в резервуаре, волны в океане);
      \item простая реализация на регулярной сетке и на GPU.
\end{itemize}

И следующие недостатки:

\begin{itemize}
      \item модель предполагает малую глубину и не описывает
            сложные 3D эффекты;
      \item трудно получить вихревую структуру потока
            внутри объёма жидкости;
      \item ограниченная физическая точность в случаях,
            выходящих за рамки предположений модели.
\end{itemize}

Для задач визуализации волн на плоскости модель мелкой
воды является очень привлекательной, но менее универсальной,
чем общий решатель уравнений Навье Стокса.

\subsection{FLIP (Fluid-Implicit-Particle)}

Метод FLIP сочетает Эйлерову сетку и лагранжевы
частицы. Сетка используется для решения уравнений
давления и несжимаемости, а частицы несут массу и
импульс. На каждом шаге:

\begin{itemize}
      \item величины переносятся с частиц на сетку;
      \item на сетке решается задача проекции скорости;
      \item обновлённая скорость интерполируется обратно
            к частицам;
\end{itemize}

Метод FLIP обладает рядом преимуществ. Во-первых, он обеспечивает меньшую численную диффузию по сравнению с чисто сеточными схемами, что позволяет лучше сохранять мелкие вихри и детали потока. Во-вторых, FLIP удачно сочетает стабильность вычислений с высокой степенью детализации, делая его подходящим для реалистичного моделирования жидкостей. То есть этот метод хорошо подходит для создания анимаций воды с брызгами и сложными поверхностными эффектами, однако требует сложной реализации, а также значительных вычислительных ресурсов и памяти для хранения как частиц, так и сеточных данных.

\subsection{MLS-MPM (Material Point Method)}

Метод MLS-MPM можно рассматривать как развитие частично\-матричных
подходов и, в частности, как дальнейшее развитие идей SPH, где
непрерывная среда представляется набором частиц. Как и в SPH, материальные
точки (частицы) хранят массу и скорость, однако в MLS-MPM они дополнительно
несут информацию о деформации, а взаимодействие между частицами
осуществляется не напрямую, а через вспомогательную эйлерову сетку.

Базовый MPM (Material Point Method) представляет материал в виде
набора материальных точек, которые сами по себе не сталкиваются
друг с другом, а обмениваются импульсом и силами через фоновую
решётку. В большинстве вариантов реализации MLS-MPM для переноса величин между частицами
и узлами сетки используется метод наименьших квадратов, что повышает точность и устойчивость схемы
при сильных деформациях.

Данный подход получил широкую известность благодаря использованию
MPM для моделирования снега в фильме Frozen студии Disney,
что наглядно демонстрирует способность метода правдоподобно
описывать широкий спектр материалов \textemdash{} от сыпучих сред до вязких
жидкостей\cite{webgpuarticle}.

Ключевыми преимуществами MLS-MPM являются:
единый формализм для жидкостей, упругих и пластичных материалов,
хорошая устойчивость при сильных деформациях и высокая визуальная
реалистичность сложных эффектов (снег, грязь, вязкие жидкости).
В то же время метод отличается существенной математической и
алгоритмической сложностью, его сложнее эффективно реализовать
в браузере, да и по отношению к цели простой демонстрации течения несжимаемой жидкости он часто
оказывается избыточным.


\subsection{Обоснование выбора метода}

С точки зрения баланса между качеством визуализации,
сложностью реализации и производительностью в браузере
наиболее подходящим является
сеточный метод класса Stable Fluids.

\begin{itemize}
      \item Алгоритм напрямую реализует упрощённые
            уравнения Навье Стокса для несжимаемой жидкости,
            что позволяет связать визуализацию с физической
            моделью.
      \item Регулярная сетка и локальные операции хорошо
            ложатся на модель вычислительных шейдеров WebGPU.
      \item Метод относительно прост для реализации,
            включает ограниченное число стадий (адвекция,
            диффузия, проекция), которые удобно разделить на
            отдельные шейдеры.
      \item Получаемые эффекты (вихревые структуры, плавные
            завихрения дыма и жидкости) визуально выразительны
            и заметны пользователю.
\end{itemize}

Модели мелкой воды и частичные методы (SPH, FLIP,
MLS \textemdash{}MPM) представляют собой интересные альтернативы
и могут быть рассмотрены как направления дальнейшего
расширения проекта. Однако в рамках данной курсовой
работе основным выбран \emph{Эйлеров сеточный метод
      Stable Fluids}, реализуемый на WebGPU.
